% Generated draft: requires mathematical and source review.
\documentclass[11pt]{article}
\usepackage[margin=25mm]{geometry}
\usepackage{fontspec}
\setmainfont{Latin Modern Roman}
\usepackage{amsmath,amssymb,mathtools,unicode-math}
\setmathfont{Latin Modern Math}
\usepackage{enumitem,needspace,xurl,hyperref}
\setlength{\parindent}{0pt}
\setlength{\parskip}{5pt}
\setlength{\emergencystretch}{3em}
\newcommand{\N}{\mathbb{N}}
\newcommand{\Z}{\mathbb{Z}}
\newcommand{\Q}{\mathbb{Q}}
\newcommand{\R}{\mathbb{R}}
\newcommand{\C}{\mathbb{C}}
\newcommand{\F}{\mathbb{F}}
\newcommand{\A}{\mathbb{A}}
\newcommand{\PP}{\mathbb{P}}
\newcommand{\E}{\mathbb{E}}
\newcommand{\eps}{\varepsilon}
\newcommand{\dd}{\,\mathrm{d}}
\begin{document}
\section*{UnsolvedMath ID: 701496}
\textbf{Source problem code:} catalogue entry
\subsection{Definitions and mathematical statement}
For a finite group \(K\), let \(\exp(K)=\operatorname{lcm}\{|x|:x\in K\}\), with \(\exp(1)=1\). If \(A\leq\operatorname{Aut}(G)\), define its common fixed-point subgroup by
\[
C_G(A)=\{g\in G:a(g)=g\text{ for every }a\in A\}.
\]
A finite Frobenius group \(FH=F\rtimes H\) has nontrivial kernel \(F\) and nontrivial complement \(H\), with \(C_F(h)=1\) for every \(h\in H\setminus\{1\}\), where \(h\) acts on \(F\) by conjugation.

Does there exist a function
\[
b:\mathbb{Z}_{\geq1}\times\mathbb{Z}_{\geq1}\longrightarrow\mathbb{Z}_{\geq1}
\]
such that, for every finite group \(G\) and every Frobenius subgroup \(FH\leq\operatorname{Aut}(G)\) with kernel \(F\), complement \(H\), and \(C_G(F)=1\), one has
\[
\exp(G)\leq b\bigl(|F|,\exp(C_G(H))\bigr)?
\]
The bound is required to depend only on the two displayed parameters.

The supplied record labels the general question open and reports affirmative results when \(F\) is cyclic, and when \(FH\cong A_4\) with \(\gcd(|G|,3)=1\), where \(A_4\) is the alternating group on four letters. These literature assessments are supplied source claims, not independently verified conclusions. The source's malformed expression for the latter group is interpreted as \(A_4\).
\subsection{Short English statement}
Must the exponent of a finite group admitting such a Frobenius group of automorphisms be bounded solely by the order of the Frobenius kernel and the exponent of the complement's fixed-point subgroup? The kernel must have no nonidentity common fixed point in \(G\). The supplied record reports two affirmative special cases and labels the general question open.
\subsection{Sources}
\begin{itemize}
\item Local UnsolvedMath registry, permanent ID \(701496\), exact source code \texttt{catalogue entry}, title ``Kourovka Notebook Problem 18.67.'' The supplied record attributes the problem to N. Yu. Makarenko, E. I. Khukhro, and P. Shumyatsky, with proposed year 2014; its supplied status is \texttt{open}.
\item \emph{The Kourovka Notebook (Unsolved Problems in Group Theory)}, original problem 18.67, 18th issue (2014). The supplied background also identifies the 21st issue edition (2026), arXiv:1401.0300v48; the supplied collection description names E. I. Khukhro and V. D. Mazurov as editors. These source details have not been independently checked.
\item E. I. Khukhro, N. Y. Makarenko, and P. Shumyatsky, \emph{Forum Math.} \textbf{26} (2014), 73--112. Supplied reference for the case of cyclic \(F\); no article title or link is supplied.
\item P. Shumyatsky, \emph{J. Algebra} \textbf{331} (2011), 482--489. Supplied reference for \(FH\cong A_4\) under the hypothesis \(\gcd(|G|,3)=1\); no article title or link is supplied.
\end{itemize}
\end{document}
